%% ============================================================
%% Chapter 1: The Latency of Evidence
%% ============================================================

\chapter{The Latency of Evidence}
\label{ch:ch01-latency-of-evidence}
\section{The Interval Between Existence and Knowledge}

It is tempting to treat discovery as the moment at which something new comes into being. An unnoticed entry in a maintenance log is read for the first time and a mechanical failure is explained; a photograph taken decades earlier is re-examined and a disputed sequence of events is settled; a sample preserved for one purpose is reanalyzed under a later technique and yields a result no one had thought to ask for. In each case it is natural to say that something was found. It is more accurate, though less natural, to say that nothing newly came into existence. The log entry, the photograph, and the sample were already there. What changed was the relation between that material and an inquiring system, not the material itself.

This chapter treats that relation as having a temporal structure of its own, distinct from the temporal structure of the event the material bears on. A trace may become materially available at one time, recoverable at a later time, warrantedly informative about some proposition at a later time still, and institutionally acted upon later than that, or it may never complete this sequence at all. What follows develops this structure precisely enough to say something sharper than that knowledge lags behind evidence: the stages a trace passes through are not independent way-stations visited in arbitrary order. They form a dependency chain in which each later stage ordinarily presupposes the one before it, even though none of the earlier stages guarantees that the later ones will ever be reached.

\section{From Trace to Judgment: Four Objects, Not Four Synonyms}

A recurring source of confusion, in ordinary usage as much as in careless formal treatments, is the habit of using a single word, \emph{evidence}, for objects that should be kept apart. An event occurs. The event leaves a trace, some physical or informational residue causally downstream of it. An interpreter may infer a proposition from that trace. An institution, community, or individual may then form a judgment that incorporates the proposition. These four objects, event, trace, proposition, and judgment, are connected but not identical:
\[
\text{event} \longrightarrow \text{trace} \longrightarrow \text{proposition} \longrightarrow \text{judgment}.
\]
A trace exists once it has been physically or informationally instantiated and persists, however briefly, in some accessible medium. A proposition is a candidate description of the event, inferable in principle from the trace. Evidence, properly speaking, is not the trace itself but the trace considered in relation to a specific proposition it bears on. A fact, depending on the sense intended, is either a true proposition or the state of affairs that makes a proposition true; what follows uses it in the first sense, as a true proposition about an event, remaining agnostic about the deeper metaphysical questions the second sense would raise.

This vocabulary supports a distinction that looser treatments of the topic collapse.

\begin{definition}[Actual and assigned support]
Let $S(e,p) \in \{0,1\}$ denote the actual, mind-independent evidential support relation: whether trace $e$ genuinely bears on proposition $p$, as a matter of causal or statistical fact, whether or not anyone has recognized it. Let $\widehat{S}_t(e,p) \in \{0,1\}$ denote the evidential status some interpreter or institution assigns to the pair $(e,p)$ at time $t$: whether $e$ is currently treated as supporting $p$.
\end{definition}

A trace may already bear evidentially on $p$ in the sense of $S(e,p)=1$ while $\widehat{S}_t(e,p)=0$, because no one has yet identified the relation. This is what it means for a trace to be evidence before anyone treats it as such, and the distinction between $S$ and $\widehat{S}$ does the load-bearing work throughout the chapter.

\section{A Dependency Structure, Not Four Independent States}

With this vocabulary in place, the chapter's four thresholds can be stated precisely. Let $X_t(e)$ hold when trace $e$ remains materially available at $t$. Let $L_t(e,p)$ hold when $e$'s bearing on proposition $p$ is warrantedly recoverable at $t$, in a sense made precise in Section~\ref{sec:legibility}. Let $A_t(e,p,c)$ hold when community or institution $c$ permits $e$ to count toward a judgment concerning $p$ at $t$. Let $C_t(e,p,c,d)$ hold when $e$'s admission by $c$ actually changes decision or state $d$ at $t$.

It is tempting to treat these four statuses as logically independent, occurring in any order. That claim does not survive scrutiny. Consequence ordinarily presupposes some operative decision process into which the material has entered, which presupposes that the material has been admitted; admissibility presupposes some intelligible claim under consideration, which presupposes legibility; and legibility presupposes a trace that exists or formerly existed.

\begin{proposition}[The evidential dependency chain]
\label{prop:dependency-chain}
The four thresholds satisfy a chain of defeasible implications running from consequence back to availability,
\[
C_t(e,p,c,d) \Rightarrow A_t(e,p,c) \Rightarrow L_t(e,p) \Rightarrow X_t(e),
\]
together with the corresponding non-implications in the forward direction,
\[
X \not\Rightarrow L, \qquad L \not\Rightarrow A, \qquad A \not\Rightarrow C.
\]
\end{proposition}

This is a more modest claim than mutual independence, and it is the claim the rest of the chapter needs. A trace can exist without being legible, be legible without being admissible, and be admissible without becoming consequential; but a trace cannot become consequential without first having been admitted, admitted without first having been legible, or legible without having existed, at least at some earlier time, in some accessible form.

What can appear to violate this ordering is retrospective attribution, in which a later interpreter judges that an old trace would have satisfied an earlier threshold had the right method or standing existed at the time: a court may recognize that a decades-old sample would have been admissible under a rule adopted only recently, or an institution may act on material before formally admitting it through its own procedures. These are real and important phenomena, but they are cases of a later re-evaluation of an earlier trace, not cases of consequence literally preceding existence.

\section{Legibility, Warrant, and Two Kinds of Recovery}
\label{sec:legibility}

The concept doing the most work in this chapter is legibility, and it is worth being careful about how it is defined. A definition stating only that some available procedure produces a proposition from a trace,
\[
R_t(e,p) :\Leftrightarrow \exists\, q \in Q_t,\ I \in \mathcal{I}_t \text{ such that } I(q,e) = p,
\]
is satisfied by any procedure that outputs $p$ when queried, whether or not the output is warranted. A confabulating interpreter, human or computational, can satisfy $R_t(e,p)$ as easily as a careful one. $R_t$ is therefore properly understood as \emph{recoverability}, not legibility: it says only that $p$ is producible, not that $p$ is supported.

Legibility requires an additional condition, a warrant relation $W_t(e,p,I)$ asserting that the interpretation was produced by a procedure reliable enough, under an explicit standard, to be trusted on pairs like $(e,p)$.

\begin{definition}[Legibility]
Legibility is recoverability plus warrant:
\[
L_t(e,p) :\Leftrightarrow R_t(e,p) \land W_t(e,p,I).
\]
\end{definition}

A further distinction is needed within legibility itself, because the scientific and debugging cases discussed below expose a category that a single undifferentiated notion of legibility conceals. Semantic legibility concerns whether a trace's immediate content can be recovered at all: whether a script can be decoded, a number read, a log line parsed. Evidential legibility concerns whether the trace's warranted bearing on a specific proposition $p$ can be established:
\[
L^s_t(e) :\Leftrightarrow \text{the content of } e \text{ can be decoded},
\]
\[
L^e_t(e,p) :\Leftrightarrow L^s_t(e) \land W_t(e,p,I) \text{ for some } I \text{ recovering } p \text{ from } e.
\]
An undeciphered inscription lacks semantic legibility outright. A fully readable laboratory notebook containing an unrecognized anomaly possesses semantic legibility while lacking evidential legibility with respect to a theory not yet formulated. Much of what a casual description of this topic would call illegibility is, on inspection, unrecognized evidential relevance rather than an inability to read the material at all, and keeping the two apart is what allows the scientific and debugging cases below to be stated precisely.

\section{Availability Is Not Application}

A second refinement concerns where latency actually comes from. It is tempting to treat the absence of a query as the sole source of delay, and while this is the most distinctive part of the present account, it becomes reductive if left as the whole explanation. A query can already belong to $Q_t$, in the sense that it is conceptually and procedurally available to some potential interpreter, without anyone actually applying it to a given trace, because investigation is costly, an archive is poorly indexed, access is restricted, incentives favor not looking, or an institution is unwilling to reconsider a settled conclusion. Membership in $Q_t$ does not imply application:
\[
q \in Q_t \not\Rightarrow \mathrm{Apply}_t(q,e).
\]

A more complete definition of practical legibility therefore requires access to the trace and actual application of an available, warranted procedure, not merely the abstract existence of such a procedure somewhere in the relevant field:
\[
L_t(e,p) :\Leftrightarrow \exists\, q \in Q_t,\ I \in \mathcal{I}_t \Big[ \mathrm{Access}_t(e) \land \mathrm{Apply}_t(I,q,e) \land I(q,e)=p \land W_t(e,p,I) \Big].
\]

This decomposition separates delay by cause rather than treating every instance of latency as a single conceptual failure to formulate the right question. \emph{Conceptual latency} occurs when $Q_t$ genuinely lacks the query. \emph{Technical latency} occurs when $Q_t$ contains the query but $\mathcal{I}_t$ lacks a reliable procedure to execute it. \emph{Archival latency} occurs when $e$ exists but $\mathrm{Access}_t(e)$ fails, because the material is poorly catalogued or physically difficult to reach. \emph{Institutional latency} occurs when access and method both exist but the institution declines to apply them, for reasons of cost, incentive, or resistance to revisiting a settled matter. \emph{Attentional latency} occurs when nothing prevents application except that no one has yet directed attention to that trace among the many competing for it. These are different mechanisms with different remedies, and treating them as a single phenomenon obscures exactly the distinctions that make the framework useful.

\section{Admissibility as Indexed Uptake}

Admissibility is never absolute, and treating it as a single global threshold overstates the uniformity across the domains this chapter and its companions in Part~III discuss. A trace may be admissible in historical scholarship, inadmissible in a courtroom, sufficient to justify an engineering rollback, and insufficient for publication in a specific scientific venue, all at the same time. Admissibility should therefore be indexed to the community or institution $c$ doing the admitting and the decision or judgment $d$ toward which the material is being admitted:
\[
A_t(e,p \mid c, d).
\]

Legal admissibility, in particular, is a rule-governed procedural category with its own doctrinal history, and it is not simply one instance of a single generic phenomenon shared uniformly with scientific acceptance or historical credibility; those are related but analytically distinct forms of what might be called uptake, the general phenomenon of some community coming to treat a proposition as counting toward a judgment. Reserving admissibility for the specific, institutionally regulated form found in law, and using uptake as the broader term, avoids overstating the uniformity of a phenomenon that, on inspection, takes recognizably different shapes in different settings. This indexing recurs in Chapter~\ref{ch:ch10-clio-architecture}'s treatment of commitment as distinct from mere admissibility.

\section{The Archive as a Reservoir of Unrealized Relations}

With $L_t$, $A_t$, and the availability and application distinctions in place, a common intuition about archives can be restated more carefully. An archive is not memory in the full sense the word usually carries, because memory implies an active and at least readily exercisable relation between a rememberer and a remembered content. An archive is closer to a reservoir of traces for which $S(e,p)=1$ may already hold for propositions $p$ that no query in $Q_t$ yet targets, or that $Q_t$ targets without $\mathrm{Access}_t(e)$ or $\mathrm{Apply}_t$ having occurred. An archive's present usefulness is measured by how much of its holdings currently satisfy $L_t$; an archive's future value is a largely independent quantity, governed by how much of its holdings could come to satisfy $L_t$ under a $Q_{t'}$, $\mathcal{I}_{t'}$, and pattern of access and application that do not yet exist. Because $Q_t$ has historically tended to expand, this future value is real even where it cannot presently be measured, and the point applies equally to formal archives and to structures that accumulate traces incidentally, such as operational logs or physical debris not yet recognized as significant.

\section{Scientific Anomalies: Semantic Legibility Without Evidential Legibility}

The history of science offers repeated instances in which a measurement was recorded accurately, preserved without controversy, and yet remained scientifically inert for years or decades. It is tempting to describe such measurements as illegible; the vocabulary developed above makes clear that they typically possess semantic legibility throughout. Researchers can read the number, state what instrument produced it, and record it correctly in a notebook. What is missing is evidential legibility with respect to a proposition that no theory available at the time supplies: no $q \in Q_{t_0}$ exists under which the measurement would count as support for or against a specific hypothesis, so it is catalogued as noise or instrument error, which is itself a judgment, made under $Q_{t_0}$, that $S(e,p)=0$ for every $p$ then under consideration.

When a later theory arrives, supplying a new $q \in Q_{t_1}$, the same number is reread under this new query, and $L^e_{t_1}(e,p)$ may now hold for a proposition $p$ that the theory has newly made askable. The measurement did not change between $t_0$ and $t_1$. Its semantic legibility did not change either; the number was always readable. What changed was $Q_t$, and specifically the arrival of a query capable of establishing $S(e,p)=1$ where before no query could establish any relation at all. Locating the change precisely in evidential legibility, rather than in an undifferentiated notion of legibility, is what makes this a defensible claim about the history of science rather than an impressionistic one.

\section{Debugging: Access, Application, and the Missing Query}

Software logs make the same structure unusually visible because the interval between availability and application can sometimes be measured directly in the timestamps involved. A fault is frequently preceded by a sequence of logged events sufficient, in combination, to diagnose its cause. These events possess semantic legibility from the moment they are written: an engineer who reads the relevant lines can parse them without difficulty. What is typically missing is not the ability to read the log but $\mathrm{Apply}_t(I,q,e)$: no one has yet directed the correlating query, drawn from a diagnostic hypothesis, at the specific lines that matter, often because the log in question was not known to be relevant until a second, more severe occurrence of the fault prompted a wider search.

The engineer's contribution at the moment of correct diagnosis is therefore usually not the production of new evidence but the construction and application of a query, drawn from an emerging hypothesis about the system's behavior, that finally establishes $L^e_t(e,p)$ for logs that had satisfied $L^s_t(e)$ all along. This has a practical implication for system design: because the query eventually needed cannot generally be anticipated at the time a log is written, logging formats should preserve enough structure and context to remain queryable under $\mathrm{Apply}_t$ for hypotheses not yet formulated, rather than only under the narrow set of queries the system's original designers had in mind, and access controls and indexing should be designed so that $\mathrm{Access}_t(e)$ does not itself become the binding constraint once a good query finally does arrive. This point recurs, in a different register, in Chapter~\ref{ch:ch12-verify-seams}'s treatment of workarounds as causal theorems.

\section{Forensic and Legal Evidence: Admissibility as Institutionally Indexed Uptake}

Legal and forensic contexts illustrate the distinction between $L_t$ and $A_t(\cdot \mid c,d)$ with unusual clarity, because admissibility in this domain is governed by explicit, codified rules rather than the gradual and informal process of scientific persuasion. A physical trace collected at a scene may achieve full evidential legibility within hours, in the sense that a trained examiner can state with warrant what the trace shows and how it bears on a specific proposition, while remaining formally inadmissible under $A_t(\cdot \mid c=\text{court}, d=\text{verdict})$ for reasons unconnected to $L_t$: a broken chain of custody, a collection method not yet validated by the relevant standards body, or a jurisdiction's specific procedural requirements.

The history of DNA evidence illustrates this without needing to be flattened into a single clean threshold crossing. DNA analysis did not simply wait, as a unified capability, to become admissible; laboratory capability, sample integrity over time, procedural access to post-conviction testing, statutory rules governing when such testing could be compelled or requested, and judicial willingness to grant relief on the basis of results all developed unevenly and on different timelines. A given biological sample might have achieved $L_t$ years before the statutory framework granting $\mathrm{Access}_t$ to post-conviction testing existed, and might have satisfied that statutory framework years before a specific court was willing to grant $C_t$ for a specific case. Each of these is a separate threshold, indexed to a separate $c$ and $d$, and the multi-threshold model developed here is better suited to this complexity than a single-interval account would be, since it represents legibility, access, admissibility, and consequence as arriving on four different timelines rather than as one delayed unveiling.

Historical testimony exhibits a related pattern with the trace replaced by an account rather than a physical residue. Testimony can achieve full evidential legibility, in the sense that a careful listener can state exactly what is being claimed and why it would bear on a disputed proposition, while remaining excluded from $A_t(\cdot \mid c,d)$ for a dominant institutional narrative, because that community is unwilling, for reasons independent of the testimony's evidential merit, to admit it toward the relevant judgment. A later shift in the composition or commitments of $c$ can produce $A_{t'}(e,p \mid c,d)$ for testimony that had satisfied $L_t(e,p)$ unchanged since $t$, which is a case of uptake changing while the underlying support relation $S(e,p)$ did not.

\section{Machine Learning: Encoding, Extractability, and Use}

The same structure recurs within machine learning systems, with the interpreter replaced by a trained model, though care is needed not to use illegibility too broadly here either. A feature can be present in raw data in several distinct senses that should not be treated as equivalent: it may be explicitly encoded in a single variable, statistically recoverable only from a combination of variables, recoverable only with information external to the dataset, or merely correlated within one particular sample and unlikely to generalize. Separately, a model may fail to use a feature because its architecture cannot represent the relevant function, because optimization fails to discover a representation that exists in principle within its hypothesis class, because the training objective does not reward extracting it, or because regularization actively suppresses it. Collapsing all of these into a single notion of illegibility conceals which mechanism is actually responsible.

A more disciplined formulation separates three distinct non-implications, mirroring the structure used throughout the chapter:
\[
\text{encoded in } D \not\Rightarrow \text{extractable by } M \not\Rightarrow \text{used for task } T.
\]
Data can encode a pattern without any particular model architecture being able to extract it; a model can be capable of extracting a pattern without its training objective ever rewarding that extraction for the task at hand. A change in architecture, training procedure, or auxiliary objective can move a feature from encoded-but-inextractable to extractable, and separately from extractable-but-unused to used, without any change to the underlying data, in direct structural parallel to a scientific theory establishing evidential legibility for an unchanged measurement or a statutory reform establishing admissibility for an unchanged sample.

\section{Latency and Reification as Complementary Errors of a Single Support Relation}

The context reification pattern examined in Chapter~\ref{ch:ch03-theatre-of-agreement}, where persona ensembles generate a false phenomenology of social confirmation, connects to the present argument, and the connection is exact once stated in terms of $S(e,p)$ and $\widehat{S}_t(e,p)$ rather than as a loose claim about mirrored admissibility thresholds. Evidential latency and context reification are not independent failures requiring independent thresholds; they are the two possible errors of a single classifier, the one implicitly performed whenever a system decides how to treat a pair $(e,p)$:
\[
\text{latency: } S(e,p)=1,\ \widehat{S}_t(e,p)=0, \qquad
\text{reification: } S(e,p)=0,\ \widehat{S}_t(e,p)=1.
\]

Latency is a false negative: material genuinely bears on a proposition, but the system's assigned status has not caught up. Reification is a false positive: material does not genuinely bear on a proposition, but the system's assigned status has run ahead of the evidence, typically because a generative process produced a fluent, internally consistent proposition that satisfied $R_t(e,p)$ without satisfying $W_t(e,p,I)$, in exactly the sense the recoverability-versus-warrant distinction of Section~\ref{sec:legibility} was built to capture. Framed this way, the two failures are not evidence for two separate admissibility thresholds needing separate tuning, but two error rates of one underlying support-assignment process, and different institutional or architectural controls, stricter warrant standards, broader query availability, more diligent application, better provenance tracking, will generally trade against each other along a single operating curve rather than requiring two unrelated interventions.

\section{Decomposing Latency}

Because the dependency chain of Proposition~\ref{prop:dependency-chain} orders four thresholds rather than treating them as a single undifferentiated delay, the overall interval between a trace's availability and its eventual consequence can be decomposed into three separate latencies, each attributable to a different mechanism:
\[
\lambda_L = t_L - t_X, \qquad \lambda_A = t_A - t_L, \qquad \lambda_C = t_C - t_A,
\]
where $t_X$, $t_L$, $t_A$, and $t_C$ are the times at which $X_t(e)$, $L_t(e,p)$, $A_t(e,p \mid c,d)$, and $C_t(e,p,c,d)$ first hold, when they hold at all. This decomposition gives the chapter's title concept genuine content rather than a single impressionistic gap. An undeciphered inscription exhibits large $\lambda_L$, driven by the absence of semantic legibility itself. A forensic technique that is scientifically well understood but not yet legally admissible exhibits small $\lambda_L$ and large $\lambda_A$. A historical finding that is fully admitted into scholarly consensus but never applied to revise any live policy or decision exhibits large, potentially unbounded $\lambda_C$. Characterizing a case by its latency profile, rather than by a single generic claim that knowledge lagged behind evidence, is where the framework's descriptive value actually lies.

\section{Preservation, Cost, and Option Value}

The claim that preservation carries option value follows directly from the framework: because $Q_t$ expands unpredictably, material that satisfies $X_t(e)$ without satisfying $L_t(e,p)$ for any currently formulable $p$ may still come to satisfy $L_{t'}(e,p')$ for some later $t'$ and some $p'$ not presently conceivable, and this possibility has genuine expected value even though it cannot be measured from the present. It would be a mistake, however, to move directly from this observation to the stronger and unqualified recommendation that apparently irrelevant residue should generally be retained. That recommendation does not follow without a theory of costs, and stating it without qualification risks reading as an argument for indiscriminate retention and surveillance, which is not this chapter's intent.

Preservation carries genuine costs that must be weighed against option value rather than overridden by it: storage and curation costs, exposure of private information, security risk, the retention of defamatory or coercive material, the burden placed on retrieval systems, the generation of misleading correlations from sheer volume, and in some cases a legal or ethical obligation to delete. There is also a tension internal to the framework itself between two goods that are easy to conflate: preserving more material increases the total reservoir of traces that might someday satisfy $L_t$, but it can simultaneously increase $\lambda_L$ for the traces that matter most, by burying them in a much larger volume of material that does not, and making $\mathrm{Access}_t(e)$ and $\mathrm{Apply}_t$ harder to achieve for any given $e$. Preservation and discoverability are not the same good, and a policy optimized purely for the first can degrade the second.

The defensible form of the argument is therefore a rebuttable principle rather than an unconditional imperative:
\[
\text{uncertain present relevance} \not\Rightarrow \text{zero expected future value}.
\]
Deletion decisions should weigh epistemic option value alongside privacy, security, cost, consent, and discoverability, rather than treating any one of these considerations, including future evidential value, as automatically dominant.

\section{Deletion as Reduction in Recoverability, Not Total Foreclosure}

A related overstatement to guard against is the claim that the asymmetry between illegibility and deletion is total, illegibility being reversible and deletion being an irreversible foreclosure of every future query. This is rhetorically effective but formally false in both directions. A deleted trace may survive in copies, caches, quotations, derived statistics, backups, human memory, or causally downstream records, and its content may in some cases be partially reconstructed from this redundancy, though the chain of provenance connecting the reconstruction to the original event may be degraded or lost even where the content itself survives, which can eliminate the evidentiary force of the material without eliminating its informational content. Conversely, an undeleted trace can become effectively unrecoverable through media decay, format obsolescence, loss of a decryption key, or loss of the contextual information needed to interpret it, so that its formal continued existence, $X_t(e)$, provides little practical difference from deletion.

The more accurate model replaces a binary distinction between existing and deleted with a continuous or graded measure of recoverability,
\[
\rho_t(e) \in [0,1],
\]
or, where a numerical value is unwarranted, a small ordered set of qualitative states such as intact, degraded, indirectly reconstructible, inaccessible, and destroyed. Preservation policy, on this view, is better understood as an attempt to manage the trajectory of $\rho_t(e)$ over time for the traces judged worth the associated costs, rather than as a binary decision to keep or discard, and deletion is better understood as a sharp downward step in $\rho_t(e)$, often but not always to its floor, rather than as an act of total metaphysical foreclosure. This graded recoverability measure anticipates the qualified non-erasure principle developed for witnessed history in Chapter~\ref{ch:ch10-clio-architecture}.

\section{Conclusion}

Existence, legibility, admissibility, and consequence are not four independent statuses a trace may acquire in any order; they form a dependency chain in which each later status ordinarily presupposes the one before it, while none of the earlier statuses guarantees that the later ones will ever be reached. Legibility itself divides into semantic legibility, whether a trace's content can be decoded at all, and evidential legibility, whether its warranted bearing on a specific proposition can be established, and the gap between these two is where most of what looks like undiscovered evidence actually resides: in scientific archives, in operational logs, and in machine-learning datasets alike, the material is frequently already readable, and the missing element is a query, a method, an act of application, or an institutional willingness to admit it, not the ability to perceive the material at all.

This structure explains why anomalies can sit fully readable in scientific archives for decades before a new theory establishes their evidential bearing, why faults are frequently diagnosable in retrospect from logs that had recorded everything necessary all along, why forensic and testimonial evidence can achieve evidential legibility years or decades before achieving admissibility within a specific institution, and why the same dataset can be extractable by one model architecture and inert to its predecessor. It also gives exact formal content to the connection between this chapter and the theatre of agreement examined in Chapter~\ref{ch:ch03-theatre-of-agreement}: latency and reification are not two independent thresholds needing separate correction, but the two error types, false negative and false positive, of a single evidential support-assignment process, and institutional or architectural interventions should be understood as trading between these error rates rather than as independently tunable dials.

The ethical conclusion that follows is more modest than an unqualified case for preserving everything, and better for being so. Because $Q_t$ expands in ways the present cannot foresee, uncertain present relevance does not imply zero expected future value, and deletion should be weighed as a genuine reduction in future recoverability rather than dismissed as costless simply because the material appears useless today. But this consideration sits alongside real costs, including the risk that indiscriminate retention degrades the very discoverability that gives preservation its point. A fact can already exist, fully formed and fully evidenced, in a form that no one now alive is equipped to recognize as such. The obligation this creates is not to preserve everything indefinitely, nor to read everything now, both of which are impossible, but to weigh the foreclosure of future queries as a real cost of deletion, alongside the equally real costs of retention, rather than treating either preservation or deletion as the obviously correct default.
